\documentclass[11pt]{article}
\usepackage[margin=0.75in]{geometry}
\usepackage[T1]{fontenc}
\usepackage[utf8]{inputenc}
\usepackage{setspace}
\usepackage{hyperref}
\hypersetup{colorlinks=true,urlcolor=blue,linkcolor=black}
\pagestyle{empty}
\setlength{\parskip}{0.42em}
\setlength{\parindent}{0pt}
\setstretch{0.96}
\begin{document}

\begin{flushleft}
Ke Cheng Hsu \\
Department of Economics \\
The Hong Kong University of Science and Technology \\
Clear Water Bay, Kowloon, Hong Kong \\
\href{mailto:khsu@connect.ust.hk}{khsu@connect.ust.hk} \\
\href{https://cliohsu1997.github.io/}{https://cliohsu1997.github.io/}
\end{flushleft}

\vspace{0.1em}
\today

\vspace{0.3em}
Hiring Committee \\
School of Liberal Arts and Sciences \\
Musashi University \\
Tokyo, Japan

\vspace{0.3em}
Dear Members of the Hiring Committee,

I am writing to apply for the full-time faculty position in Economics (Microeconomics and Applied Microeconomics) at Musashi University (JREC-IN D126070158). I am a Ph.D.\ candidate in Economics at the Hong Kong University of Science and Technology (expected June 2027). I am particularly drawn to this post because of its central teaching role in the School of Liberal Arts and Sciences Parallel Degree Programme with the University of London. Undergraduate microeconomics and industrial economics are at the core of that role and align closely with my field in empirical industrial organization and applied microeconomics, including research on firm and consumer behavior in product markets.

At HKUST I teach and assist in undergraduate and postgraduate economics courses in English, including microeconomic theory, information economics, and macroeconomics. In Math Camp for the MSc in Economics, I teach linear algebra, statistics, dynamic programming, real analysis, and optimization; further details are on my website. I am prepared to teach the courses listed for this appointment in English, especially \emph{Industrial Economics} (EC3099), \emph{Microeconomics} (EC2066), and the Mathematics~1 / Statistics~1 sequence. I would also be glad to contribute to SLAS administration, entrance examinations, and student supervision.

I enclose two working papers with English summaries. My job-market paper, \emph{Designing Return Policy in E-Commerce: A Trade-off between Purchase and Return Efficiency}, asks how return policy balances purchase and return efficiency. A stricter policy curbs excessive returns, but it can also raise prices and exclude valuable purchases. Using Taobao apparel data and a structural model of purchase, insurance, and return decisions, I find that the preferred policy depends on margin and scrap value. I propose a partial return-penalty discount that raises welfare relative to existing return-shipping insurance for most categories. My second paper, \emph{Product Design with Heterogeneous Search Costs} (with Pak Hung Au), studies how firms choose product design when consumers differ in search costs. Consumers with low search costs search for niche products, while consumers with high search costs search for standardized ones. Lower search costs need not lower prices once the mix of search modes adjusts.

Two further papers are in progress. With Kohei Kawaguchi and Kei Kawai, we study resale price maintenance. It can encourage retailers to hold more inventory, but when demand is low the leftover stock cannot be cleared by cutting prices and is returned at an environmental cost. With Ce Bian and Liang Zhao, we study pass-through and tax incidence when products can be returned. A tax can raise return rates and, in some cases, produce negative pass-through.

Please find enclosed my CV, publication list, and contact details for two academic referees. Thank you for your consideration.

\vspace{0.25em}
Sincerely,

\vspace{0.65em}
Ke Cheng Hsu \\
Ph.D.\ Candidate, Department of Economics \\
The Hong Kong University of Science and Technology

\end{document}
